\section{Die Nordamerikanische Stadt}
\begin{myexercise}[Arena zur Amerikanisierung der Städte]
Überlegen Sie sich während dem \href{https://www.srf.ch/play/tv/kulturplatz/video/hochhaeuser-das-leben-in-der-vertikale?urn=urn:srf:video:a98ccb20-f5ad-46ea-8203-2688102543d4}{Dokumentarfilm von SRF}, welche Argumente aus Sicht folgender Personen für oder gegen den Bau weiterer Hochhäuser in der Schweiz sprechen.
\begin{itemize}
\item Mutter / Vater mehrerer Kinder
\item Immobilien-InvestorIn
\item StadtplanerIn
\item UmweltaktivistIn
\item Landwirt(in)
\item Urbane(r) Mode-Unternehmer(in)
\end{itemize}
Wir werden basierend auf dem Film eine Debatte zum Thema ``Manhattisierung der Schweiz'' durchführen.
\end{myexercise}

\begin{mychallenge}[Weitere Unterlagen]
\begin{itemize}
\item \href{https://www.srf.ch/kultur/gesellschaft-religion/verdichtung-in-staedten-wie-kann-zuerich-platz-fuer-zusaetzliche-70-000-menschen-schaffen}{SRF-Beitrag zur Verdichtung}
\item \href{}{...}
\end{itemize}
\end{mychallenge}


\begin{mychallenge}[{Artikel zu geplanten Hochhäusern in Zürich (25.11.2024 NZZ)}]
\begin{figure}[H]
\centering
\colorbox{white}{\includegraphics[trim={1cm 17cm 1cm 1cm},clip,width=\linewidth]{"Private/Appendix/2024-11-25 Hochhäuser Zürich"}}
\end{figure}
\end{mychallenge}


\begin{mychallenge}[{Essay über Hochhäuser in Zürich (04.09.2023 NZZ Folio)}]
\foreach \pagenum in {1,2} {
    \begin{figure}[H]
    \centering
    \colorbox{white}{\includegraphics[width=\linewidth, page=\pagenum]{"Private/Appendix/2023-09-04 NZZ Folio Hochhäuser Zürich"}}
    \end{figure}
}
\end{mychallenge}


\newpage
\section{Lernziele}
\begin{todolist}
\item (Vorlektion) Ich kann den Aufbau \textbf{nordamerikanischer Städte} sowie typische Segregationsformen anhand eines räumlichen Modells beschreiben, indem ich die Begriffe \textit{CBD (Central Business District)}, \textit{Schachbrett-Grundriss} (\textit{Gridiron Pattern}), \textit{Übergangsbereich}, \textit{Auto-Infrastruktur}, \textit{Umland}, \textit{Edge City} sowie \textit{Gated Community} räumlich korrekt einordne.
\item Ich kann Tendenzen der Amerikanisierung in der Schweiz erkennen und verorten. Ich kann benennen, wie amerikanisch diese Tendenzen wirklich sind.
\item Ich kann die Hauptargumente, welche für und gegen die Verdichtung von Städten sprechen, insbesonderen hinsichtlich dem Bau neuer Wolkenkratzer, aus der Perspektive verschiedener Bevölkerungsgruppen (Familien, Investoren, Stadtplaner, Umweltaktivisten, Landwirten etc.) erläutern.
\end{todolist}

